\begin{abstract}
Segmentation models deployed on a new domain can be improved with a few labelled cases, but an update that helps on average can still silently harm individual patients.
We make self-improvement \emph{certifiable}: an update---to a refinement agent or to the segmentation model itself---is deployed only if a Learn-then-Test procedure bounds the fraction of patients it harms (a lost lesion or reduced Dice, relative to the deployed system) below a user-chosen tolerance $\delta$ with probability $1-\alpha$; otherwise the system rolls back.
We study an FDG PET/CT tumour-segmentation system transferred to an unseen tracer (PSMA, 377 patients), with disjoint adaptation, certification and locked test sets under a pre-registered protocol.
Zero-shot, the system is safe but poor (Dice 0.11).
Certified updates to the refinement agent help only modestly ($+0.009$ Dice, 0 of 125 test patients harmed).
Certified fine-tuning of the segmenter on 126 patients raises Dice from 0.11 to 0.55 and cuts false-positive volume by 98\% on the locked test; certified at $\delta=10\%$, it harms 2.4\% of test patients (95\% upper bound 6.1\%), while at our pre-registered $\delta=5\%$ the certificate refused the switch.
On the same patients, naive self-improvement harms 11 patients and loses 20 lesion detections, and pseudo-label self-training harms 4 without gain.
Across 1,000 resampled splits and three backbones, naive selection breaks a 5\% harm promise in at least 98\% of splits and certified selection in at most 0.7\%.
The tolerance $\delta$ thus becomes a dial that trades benefit for risk, with a guarantee at every setting.
\end{abstract}
